\section{Figure and table specifications}\label{sec:figures}
Four principal figures connect the first-contact information boundary, persistent memory, and compositional prediction. Separate safety, confidence, and probing figures connect the theoretical results to their empirical diagnostics. Numerical entries are unmeasured and denoted \placeholder{}.

\paragraph{Figure 1: why the first contact is different.}
The same robot encounters two recurring surfaces and loads. A shared timeline reads \emph{cue $\rightarrow$ committed action / first contact $\rightarrow$ prediction error}. In the upper branch, an episode-reset adapter can update only after the consequential action. In the lower branch, a cue selects a distribution over existing residual memories before the action. A compact inset connects seen tuples through the shared factor kernel and initializes a residual memory for an unseen tuple; a separate interaction term retains uncertainty. Unknown conditions carry broad predictions. After contact, filtered probabilities update the context belief, while the saved pre-error probabilities weight the regression update. The caption identifies the information limit: pre-contact retrieval requires prior experience or reusable structure.

\paragraph{Figure 2: the first-contact gap and its information dependence.}
Panel (a) plots common-action first-contact error against cue reliability $\rho$, including the chance point $1/K$. Episode-reset methods coincide with their initialization in the enforced no-feedback protocol; full \method{}, history-only inference, a cue-conditioned persistent comparator, and oracle retrieval quantify where additional cue information helps. Panel (b) shows task success and the true-model decision gap on the same finite candidate set. Panel (c) compares cumulative conditional mean error with the measured posterior log-loss bound; the finite-alphabet oracle entropy diagnostic is labeled separately. A recurrence inset separates cold-start and return episodes. Confidence intervals use independent streams. The caption distinguishes true mean error from noisy outcome error, and identifies each comparison's information access.

\paragraph{Figure 3: memory changes the mechanism of adaptation.}
Panel (a) shows recovery time versus occurrence, with zero-step recoveries and censored failures visible. A controlled two-context inset compares recognition time with initial log odds and records overshoot under fixed diagnostic observations. Panel (b) shows signed $R(s)$ and $P(s)$ over time, including the pre-contact point, for first, second, and later occurrences; first-exposure selection/creation is labeled separately. Their sum $G(s)$ appears as a line. Panel (c) compares full \method{} with frozen-memory retrieval on return episodes, with shared prior centers also frozen. Supporting interference curves report immediate A-return error versus intervening B exposure under inferred and oracle memory selection.

\paragraph{Figure 4: composition before the first-ever contact.}
Panel (a) displays the $3\times3$ surface--load grid and the corresponding bipartite observation graph, marking six observed and three held-out cells; a disconnected control contains a held-out cross-component pair. Panel (b) reports first-contact error for each held-out cell from independent copies of the same pretest checkpoint. The comparisons are independent whole-context memories, additive sharing with $\kappa_\times=0$, full sharing with $\kappa_\times>0$, a continuous-context baseline, and oracle factor selection. Panel (c) plots error against the controlled interaction strength $\lambda$, alongside the precomputed model conditional variance. Prior-matched replications compare variance with average squared mean error; noisy transition error is separate. An unseen interaction retains its prior uncertainty, and the primary endpoint precedes every update from the held-out combination.

\paragraph{Figure 5: safety margins are also remembered.}
The safety figure plots intervention and task cost against measured violation rate across fixed operating points, with the time course of conservatism when a familiar condition returns. Comparisons include global, continually shared, all-memory, posterior-selected per-memory, and oracle-context filters. A cumulative-violation panel includes the bound $\delta_s T+K_T(\bar S+2\gamma)/\gamma$ only on runs satisfying its bounded-score, exact-filter, and safe-fallback assumptions. Misleading-cue intervals, slot creation, and fallback events are visible. Latent and physical violations have separate traces and uncertainty estimates. This figure is ancillary to the four principal figures.

\paragraph{Figure 6: information matters through the decision.}
Panel (a) contrasts contexts that produce different probe outcomes but share an optimal strike with contexts whose optimal strikes differ. The former have zero gross decision value despite positive mutual information in the pure-information setting. Panel (b) shows the success--probe-cost frontier for no probing, information maximization, value-of-information selection, and the linearized $c$-optimal rule at matched candidate sets and budgets. Panel (c) compares estimated and realized value while varying posterior masses, decision cost gaps, and probe informativeness. Physical-nudge trials are separate from the fixed-state diagnostic, and include motion-induced cost and estimation uncertainty.

\paragraph{Main table: task-level relevance.}
The table spans controlled simulation, adapted world-model benchmarks, and hardware. Rows group representative methods by retained information. Columns report first-contact prediction error, task success, contact-critical violations where applicable, per-decision runtime, and added memory. Return and novel-combination conditions occupy separate blocks. Unmeasured entries are \placeholder{}; unavailable comparisons are marked \emph{not evaluated}. The table provides task-level context for the controlled figures.

\paragraph{Confidence figure: calibration under uncertain assignments.}
Panel (a) compares mean bias and confidence coverage for predictive, soft filtered, hard filtered, and oracle assignment across context separation and noise. Oracle radii with valid contamination and prior terms are distinct from plug-in empirical intervals. Panel (b) shows interval width and error versus the contaminated weight fraction; panel (c) reports the fraction of streams with any failure on the finite query/time grid. The fixed-gating scalar diagnostic and adaptive-filter runs occupy separate panels. Coverage refers to the true residual mean; prediction intervals for noisy observations are reported separately. Contact phase, noise underestimation, and prior misspecification are identified explicitly.

\paragraph{Supporting figures and protocol table.}
Supporting figures contain cue-conflict posterior trajectories, paired false-split/missed-change heat maps, full interference curves, cue and posterior calibration, and correction-rank/runtime scaling. The long deployment has an aligned timeline covering every episode, with context switches, returns, day boundaries, and failures marked. A protocol table records each method's cue access, persistent state, update budget, and oracle information. Additional sensitivity studies cover memory capacity, near-duplicate contexts, smooth drift, cue granularity mismatch, and delayed identifiability.
